\section{Related work}

The comparison with existing work turns on three distinctions: the population quantity being estimated, the information supplied by the statistical experiment, and the class over which performance is uniform. Our target is the scalar optimized value in \cref{obj:def:observed-value}, with its causal interpretation established in \cref{obj:prop:identification}. The \leanref{S-1}{finite observational experiment} has categorical covariates, binary treatments, and binary outcomes. Its sample size \leanref{sym:n}{\ensuremath{n}}, number of covariate cells \leanref{sym:d}{\ensuremath{d}}, and public overlap floor \leanref{sym:\epsilon}{\ensuremath{\epsilon}} jointly index the risk comparison. Under \cref{obj:ass:iid-sampling,obj:ass:shrinking-overlap}, the uniformity domain is the observed-law class in \cref{obj:def:observed-class}. Within that class, \cref{obj:thm:matched-frontier,obj:thm:consistency-threshold} characterize squared-error risk and the conditions for uniform consistency as the number of cells grows and overlap shrinks.

The closest estimand-specific literature studies differentiability and inference for the mean outcome under an optimal treatment rule. \citet[Theorem~1 and equation~(3)]{LuedtkeVanderLaan2016OptimalValue} characterize pathwise differentiability under strictly positive propensities, bounded outcomes, and finite variance of the candidate gradient. Their condition permits treatment-effect ties when both conditional outcome variances vanish on the tie set. Their online estimator also supports asymptotic inference at exceptional laws under conditions controlling standardized increments, variance estimation, the nuisance remainder, and the value loss of estimated rules \citep[Theorem~2, conditions~(C1)--(C5)]{LuedtkeVanderLaan2016OptimalValue}. The efficiency conclusion in their Corollary~3 adds fixed positivity, boundedness, convergence, and variance conditions. These results describe inference under specified asymptotic conditions at the underlying law; the present comparison describes worst-case squared-error risk over an entire finite-cell class indexed by sample size and overlap.

Subsample aggregation provides another approach to inference for the same optimized-value estimand. \citet[Theorems~2--3]{Shi2020} combine treatment-rule estimation on subsamples with evaluation on complementary observations. Their analysis assumes fixed positivity and an expected value-loss rate with exponent greater than one half, together with compatible subsample growth \citep[conditions~(A3)--(A4)]{Shi2020}. For a rule obtained from an estimated treatment contrast, their margin and contrast-estimation conditions provide sufficient conditions for this value-loss requirement \citep[conditions~(A5)--(A6)]{Shi2020}. The feasible inference result additionally imposes variance and nuisance-estimation conditions \citep[Theorem~3 and equations~(9)--(11)]{Shi2020}. This approach explains how aggregation can accommodate nonunique optimal rules while supporting asymptotic inference. Our finite-cell risk characterization instead takes the uniformity domain directly from the observed-law and overlap restrictions.

Softmax smoothing connects optimal-value inference to estimation of a differentiable surrogate. \citet[Theorem~3.3 and Corollary~3.4]{WhitehouseChenAusternSyrgkanis2025} establish asymptotic linearity and studentized inference under strong positivity, controlled density of positive suboptimality gaps near zero, suitable growth of the inverse smoothing scale, and nuisance-estimation conditions \citep[Assumptions~1--2]{WhitehouseChenAusternSyrgkanis2025}. Their density condition permits an atom at zero, thereby accommodating ties. The theorem balances smoothing bias against a smoothing-dependent mean-square rate for the outcome regression and a product-rate condition for the regression and its representer; boundedness and positive limiting variance also enter the inference guarantee. The corresponding distinction here is the uniformity domain: \cref{obj:thm:matched-frontier} evaluates estimation risk across the finite-cell class, including laws with ties and arbitrarily small positive treatment contrasts, as the public overlap floor varies.

The nearest comparison in the discrete causal experiment is \citet{ZengBalakrishnanHanKennedy2024Discrete}. Their June 2026 revision studies arm means and average treatment effects with many discrete covariate categories. For a treated-arm mean, their Theorem~1 gives a uniform mean-squared-error upper bound of the form
\[
\frac{(1-\epsilon)^2}{\epsilon^2}\frac{d^2}{n^2}
+\frac{C}{n\epsilon},
\]
where \(C\) is an absolute constant \citep[Theorem~1]{ZengBalakrishnanHanKennedy2024Discrete}. Their minimax lower bound has order
\[
\left(\frac{d}{n\log n}\right)^2+\frac{1}{n}
\]
in the regime \(d\lesssim n\log n\), with a comparison constant depending on overlap \citep[Theorem~2]{ZengBalakrishnanHanKennedy2024Discrete}. These results expose the interaction between sparse cell counts and propensity denominators for a target that is linear in the arm-specific outcome means. The optimized value adds a maximum within each cell. Our comparison retains the overlap dependence in both sides of the risk characterization for this nonlinear target.

The approximation architecture draws on the literature on estimating functionals of discrete distributions. \citet{JiaoVenkatHanWeissman2015Functionals} combine polynomial approximation in nonsmooth regions with unbiased estimation of polynomial terms and bias correction in smooth regions. Their entropy characterization applies when the sample size is at least a constant multiple of alphabet size divided by its logarithm, and separates approximation bias from a variance contribution \citep[Theorem~1]{JiaoVenkatHanWeissman2015Functionals}. Thus the logarithmic improvement has an estimand-specific risk expression and a stated sampling regime. In our experiment, the localized approximation is applied to the armwise extension in \cref{obj:def:armwise-extension}, and the attaining construction estimates its polynomial contributions through factorial statistics \cref{obj:def:armwise-estimator}. The propensity denominators make the overlap floor part of the approximation and statistical analysis.

The \(L_1\)-distance results of \citet{JiaoHanWeissman2018L1} provide a particularly close comparison because absolute-value nonsmoothness is closely related to the maximum of two arm means. They treat both estimation relative to a known reference distribution and estimation from samples from two unknown distributions. For alphabet size \(d\), their worst-case squared-error rate \(d/(n\log n)\) applies in the regime \(n\gtrsim d/\log d\) and \(\log n\lesssim\log d\), with the precise distribution-specific qualifications stated for the known-reference problem \citep[Theorems~2--4 and~6, Remark~1]{JiaoHanWeissman2018L1}. These regime restrictions matter when comparing logarithmic improvements across estimands. Here, weighted approximation and priors with prescribed common mean and matching moments connect the scalar approximation problem to the causal experiment \cref{obj:def:weighted-error,obj:def:constrained-prior-separation,obj:thm:weighted-separation-and-frontier}.

Policy learning evaluates a different statistical decision. \citet[Assumption~2.1 and Theorems~2.1--2.2]{Kitagawa2018} establish uniform welfare-regret bounds for empirical welfare maximization over a policy class, with class complexity and overlap entering the comparison. \citet[Theorem~1]{Athey2021} use estimated doubly robust scores to obtain asymptotic regret guarantees under nuisance-estimation, policy-class, and outcome conditions. \citet{JinRenYangWang2025Pessimism} construct pessimistic policy choices using known behavior propensities, with guarantees governed by coverage of an optimal policy and the complexity of the candidate class, including settings with adaptive data collection. In these problems, the statistical output is a treatment rule, and loss measures its welfare shortfall relative to a policy-class optimum. The risk in \cref{obj:def:minimax-risk} instead measures squared error in a scalar population value optimized over all cellwise treatment choices.

Specified-policy evaluation offers a further useful distinction. \citet{Ma2022} characterize minimax performance for evaluating a given target policy in a multi-armed bandit experiment, separating known, unknown, and partially known behavior policies. Their partial-knowledge analysis also uses polynomial approximation when a public lower bound on behavior probabilities is available \citep[Section~3.3]{Ma2022}. This supplies a related treatment of limited sampling probabilities, while the present target incorporates optimization over the conditional outcome means.

The closest same-target predecessor is the fixed-overlap result of \citet[Theorem~3]{CausalSmithFixedOverlap2026}. For every fixed \(\epsilon\in(0,1/2)\), it proves the rate \(\min\{1,d/[n\log(ed)]\}\), with comparison constants that may depend on \(\epsilon\), using the same Jackson--factorial architecture. The present result identifies the joint rate for \(0<\epsilon\le1/2\) when that floor may shrink: \(\min\{1,d/[n\epsilon\log(ed)]\}\) with universal comparison constants, and the corresponding consistency threshold in \cref{obj:thm:consistency-threshold}. The additional argument controls the maximum through separate armwise sensitivities, constructs weighted moment-matched priors with prescribed common mean, and compares the full observed likelihood at boundary propensities. Thus the fixed-overlap specialization is recovered while the dependence on overlap is made uniform and explicit.
